\chapter{2004年湖北卷（理）}
\section{单选题}
\begin{problem}
与直线 \(2x - y + 4 = 0\) 平行的抛物线 \(y = x^2\) 的切线方程是（\(\quad\)）

\choices
    {\(2x - y + 3 = 0\)}
    {\(2x - y - 3 = 0\)}
    {\(2x - y + 1 = 0\)}
    {\(2x - y - 1 = 0\)}
\end{problem}

\begin{answer}
D
\end{answer}

\begin{solution}
设切点为 \((x_0,x_0^2)\)，则抛物线在该点处的切线为 \(y=2x_0x-x_0^2\)，其斜率为 \(2x_0\). 因为切线与直线 \(2x-y+4=0\) 平行，所以 \(2x_0=2\)，即 \(x_0=1\). 故切线方程为 \(y=2x-1\)，即 \(2x-y-1=0\)，故选 D.
\end{solution}

\begin{problem}
复数 \(\frac{(-1 + \sqrt{3}\i)^5}{1 + \sqrt{3}\i}\) 的值是（\(\quad\)）

\choices
    {\(-16\)}
    {\(16\)}
    {\(-\frac{1}{4}\)}
    {\(\frac{1}{4} - \frac{\sqrt{3}}{4}\i\)}
\end{problem}

\begin{answer}
A
\end{answer}

\begin{solution}
由 \(-1+\sqrt{3}\i=2\left(\cos\frac{2\pi}{3}+\i\sin\frac{2\pi}{3}\right)\)，得
\( (-1+\sqrt{3}\i)^5=32\left(\cos\frac{4\pi}{3}+\i\sin\frac{4\pi}{3}\right)=-16-16\sqrt{3}\i\). 因此
\(\displaystyle\frac{(-1+\sqrt{3}\i)^5}{1+\sqrt{3}\i}=\frac{(-16-16\sqrt{3}\i)(1-\sqrt{3}\i)}{(1+\sqrt{3}\i)(1-\sqrt{3}\i)}=-16\)，故选 A.
\end{solution}

\begin{problem}
已知 \( f\left(\frac{1-x}{1+x}\right)=\frac{1-x^{2}}{1+x^{2}} \)，则 \( f(x) \) 的解析式可取为（\(\quad\)）

\choices
    {\(\frac{x}{1 + x^2}\)}
    {\(-\frac{2x}{1 + x^2}\)}
    {\(\frac{2x}{1 + x^2}\)}
    {\(-\frac{x}{1 + x^2}\)}
\end{problem}

\begin{answer}
C
\end{answer}

\begin{solution}
令 \(t=\dfrac{1-x}{1+x}\)，则 \(x=\dfrac{1-t}{1+t}\). 代入右端可得
\(\displaystyle\frac{1-x^2}{1+x^2}=\frac{(1+t)^2-(1-t)^2}{(1+t)^2+(1-t)^2}=\frac{2t}{1+t^2}\). 所以 \(f(t)=\dfrac{2t}{1+t^2}\)，将自变量记为 \(x\)，即 \(f(x)=\dfrac{2x}{1+x^2}\)，故选 C.
\end{solution}

\begin{problem}
已知 \(\bs{a}\)，\(\bs{b}\)，\(\bs{c}\) 为非零的平面向量. 甲： \(\bs{a} \cdot \bs{b} = \bs{a} \cdot \bs{c}\)，乙： \(\bs{b} = \bs{c}\)，则（\(\quad\)）

\choices
    {甲是乙的充分条件但不是必要条件}
    {甲是乙的必要条件但不是充分条件}
    {甲是乙的充要条件}
    {甲既不是乙的充分条件也不是乙的必要条件}
\end{problem}

\begin{answer}
B
\end{answer}

\begin{solution}
若乙成立，即 \(\bs{b}=\bs{c}\)，则必有 \(\bs{a}\cdot\bs{b}=\bs{a}\cdot\bs{c}\)，所以甲是乙的必要条件. 反之，甲不一定推出乙，例如取 \(\bs{a}=(1,0)\)，\(\bs{b}=(0,1)\)，\(\bs{c}=(0,-1)\)，则三向量均非零且 \(\bs{a}\cdot\bs{b}=\bs{a}\cdot\bs{c}=0\)，但 \(\bs{b}\ne\bs{c}\). 因此甲不是乙的充分条件，故选 B.
\end{solution}

\begin{problem}
若 \(\frac{1}{a} < \frac{1}{b} < 0\)，则下列不等式：

\begin{circlelist}
\item \(a + b < ab\)；
\item \(|a| > |b|\)；
\item \(a < b\)；
\item \(\frac{b}{a} + \frac{a}{b} > 2\).
\end{circlelist}

其中正确的不等式有（\(\quad\)）

\choices
    {\(1\) 个}
    {\(2\) 个}
    {\(3\) 个}
    {\(4\) 个}
\end{problem}

\begin{answer}
B
\end{answer}

\begin{solution}
由 \(\dfrac{1}{a}<\dfrac{1}{b}<0\) 知 \(a<0\)，\(b<0\). 又因 \(ab>0\)，不等式乘以 \(ab\) 后得 \(b<a\)，所以 \(|a|<|b|\)，不等式 2、3 均不成立. 由于 \(a+b<0<ab\)，不等式 1 成立. 令 \(t=\dfrac{b}{a}>0\)，由 \(a\ne b\) 知 \(t\ne1\)，于是由基本不等式 \(t+\dfrac1t>2\)，故不等式 4 也成立. 因此共有两个正确的不等式，故选 B.
\end{solution}

\begin{problem}
已知椭圆 \(\frac{x^2}{16} + \frac{y^2}{9} = 1\) 的左，右焦点分别为 \(F_1\)，\(F_2\)，点 \(P\) 在椭圆上，若 \(P\)，\(F_1\)，\(F_2\) 是一个直角三角形中的三个顶点，则点 \(P\) 到 \(x\) 轴的距离为（\(\quad\)）

\choices
    {\(\frac{9}{5}\)}
    {\(3\)}
    {\(\frac{9\sqrt{7}}{7}\)}
    {\(\frac{9}{4}\)}
\end{problem}

\begin{answer}
D
\end{answer}

\begin{solution}
椭圆的焦距参数满足 \(c^2=4^2-3^2=7\)，所以 \(F_1=(-\sqrt7,0)\)，\(F_2=(\sqrt7,0)\). 若直角在点 \(P\)，则 \(PF_1^2+PF_2^2=F_1F_2^2=28\)，而 \(PF_1+PF_2=8\). 这将推出 \(PF_1PF_2=18\)，与 \((PF_1+PF_2)^2\geqslant4PF_1PF_2\) 矛盾，故直角只能在 \(F_1\) 或 \(F_2\). 此时 \(PF_i\perp F_1F_2\)，从而 \(P\) 的横坐标为 \(\pm\sqrt7\). 代入椭圆方程得 \(\dfrac{7}{16}+\dfrac{y^2}{9}=1\)，即 \(y^2=\dfrac{81}{16}\)，所以点 \(P\) 到 \(x\) 轴的距离为 \(\dfrac94\)，故选 D.
\end{solution}

\begin{problem}
函数 \( f(x) = a^x + \log_a(x + 1) \) 在 \([0, 1]\) 上的最大值与最小值之和为 \(a\)，则 \(a\) 的值为（\(\quad\)）

\choices
    {\(\frac{1}{4}\)}
    {\(\frac{1}{2}\)}
    {\(2\)}
    {\(4\)}
\end{problem}

\begin{answer}
B
\end{answer}

\begin{solution}
在所给选项中 \(a\ne1\). 函数导数为 \(f'(x)=a^x\ln a+\dfrac{1}{(x+1)\ln a}\). 当 \(0<a<1\) 时，\(f'(x)<0\)，函数递减，最大值为 \(f(0)=1\)，最小值为 \(f(1)=a+\log_a2\)；当 \(a>1\) 时，\(f'(x)>0\)，函数递增，最小值为 \(f(0)=1\)，最大值为 \(f(1)=a+\log_a2\). 逐一代入四个选项，最大值与最小值之和分别为 \(\dfrac34\)、\(\dfrac12\)、\(4\)、\(\dfrac{11}{2}\)，只有 \(a=\dfrac12\) 时该和等于 \(a\)，故选 B.
\end{solution}

\begin{problem}
设数列 \(\{a_{n}\}\) 的前 \(n\) 项和 \(S_{n} = a\left[2 - \left(\frac{1}{2}\right)^{n - 1}\right] - b\left[2 - (n + 1)\left(\frac{1}{2}\right)^{n - 1}\right]\) （\(n = 1\)，\(2\)，\(\dots\)），其中 \(a\)，\(b\) 是非零常数. 则存在数列 \(\{x_{n}\}\)，\(\{y_{n}\}\) 使得（\(\quad\)）

\choices
    {\(a_{n} = x_{n} + y_{n}\)，其中 \(\{x_{n}\}\) 为等差数列，\(\{y_{n}\}\) 为等比数列}
    {\(a_{n} = x_{n} + y_{n}\)，其中 \(\{x_{n}\}\) 和 \(\{y_{n}\}\) 都为等差数列}
    {\(a_{n} = x_{n}\cdot y_{n}\)，其中 \(\{x_{n}\}\) 为等差数列，\(\{y_{n}\}\) 为等比数列}
    {\(a_{n} = x_{n}\cdot y_{n}\)，其中 \(\{x_{n}\}\) 和 \(\{y_{n}\}\) 都为等比数列}
\end{problem}

\begin{answer}
C
\end{answer}

\begin{solution}
将已知的前 \(n\) 项和整理为 \(S_n=2(a-b)+\bigl[b(n+1)-a\bigr]\left(\dfrac12\right)^{n-1}\). 由 \(a_1=S_1=a\)，以及 \(a_n=S_n-S_{n-1}\)（\(n\geqslant2\)），可得对一切 \(n\geqslant1\)，\(a_n=(a+b-bn)\left(\dfrac12\right)^{n-1}\). 其中 \(x_n=a+b-bn\) 是等差数列，\(y_n=\left(\dfrac12\right)^{n-1}\) 是等比数列，因此 \(a_n=x_ny_n\)，故选 C.
\end{solution}

\begin{problem}
函数 \( f(x) = ax^3 + x + 1 \) 有极值的充要条件是（\(\quad\)）

\choices
    {\(a > 0\)}
    {\(a \geqslant 0\)}
    {\(a < 0\)}
    {\(a \leqslant 0\)}
\end{problem}

\begin{answer}
C
\end{answer}

\begin{solution}
\(f'(x)=3ax^2+1\). 当 \(a\geqslant0\) 时，\(f'(x)>0\) 对任意实数 \(x\) 成立，函数严格递增，没有极值. 当 \(a<0\) 时，方程 \(f'(x)=0\) 有两个根 \(x=\pm\dfrac{1}{\sqrt{-3a}}\)，且导数符号依次为负、正、负，函数在这两个点分别取得极小值和极大值. 因此函数有极值的充要条件为 \(a<0\)，故选 C.
\end{solution}

\begin{problem}
设集合 \(P = \{m \mid -1 < m < 0\}\)，\(Q = \{m \in \R \mid mx^2 + 4mx - 4 < 0 \text{ 对任意实数 } x \text{ 恒成立}\}\)，则下列关系中成立的是（\(\quad\)）

\choices
    {\(P \subsetneq Q\)}
    {\(Q \subsetneq P\)}
    {\(P = Q\)}
    {\(P \cap Q = \varnothing\)}
\end{problem}

\begin{answer}
A
\end{answer}

\begin{solution}
若 \(m>0\)，二次式 \(mx^2+4mx-4\) 在 \(|x|\) 足够大时为正，不可能对任意实数 \(x\) 小于零. 若 \(m=0\)，该式恒为 \(-4<0\)，所以 \(0\in Q\). 若 \(m<0\)，二次式开口向下，其最大值在 \(x=-2\) 处，最大值为 \(-4m-4\). 要使它恒小于零，需 \(-4m-4<0\)，即 \(m>-1\). 因此 \(Q=(-1,0]\)，而 \(P=(-1,0)\)，故 \(P\subsetneq Q\)，选 A.
\end{solution}

\begin{problem}
已知平面 \(\alpha\) 与 \(\beta\) 所成的二面角为 \(80^{\circ}\)，\(P\) 为 \(\alpha\)，\(\beta\) 外一定点，过点 \(P\) 的一条直线与 \(\alpha\)，\(\beta\) 所成的角都是 \(30^{\circ}\)，则这样的直线有且仅有（\(\quad\)）

\choices
    {\(1\) 条}
    {\(2\) 条}
    {\(3\) 条}
    {\(4\) 条}
\end{problem}

\begin{answer}
D
\end{answer}

\begin{solution}
取两平面的单位法向量 \(\bs{n}_\alpha\)，\(\bs{n}_\beta\)，使它们的夹角为 \(80^{\circ}\). 直线方向的单位向量记为 \(\bs{u}\). 直线与每个平面成 \(30^{\circ}\)，等价于 \(|\bs{u}\cdot\bs{n}_\alpha|=|\bs{u}\cdot\bs{n}_\beta|=\cos60^{\circ}=\dfrac12\). 给两个内积分别指定符号时，\(\bs{u}\) 在两个法向量所张成平面内的投影唯一；同号时投影长度的平方为 \(\dfrac{1}{4\cos^2 40^{\circ}}<1\)，异号时为 \(\dfrac{1}{4\sin^2 40^{\circ}}<1\)，所以垂直于该平面的分量均有两个相反方向. 共有四种符号组合，但一对相反符号给出同一条无向直线，故得到 \(4\) 条过点 \(P\) 的直线，选 D.
\end{solution}

\begin{problem}
设 \( y = f(t) \) 是某港口水的深度 \( y \) （米） 关于时间 \( t \) （时） 的函数. 其中 \( 0 \leqslant t \leqslant 24 \). 下表是该港口某一天从 \(0\) 时至 \(24\) 时记录的时间 \( t \) 与水深 \( y \) 的关系：

\begin{center}
\begin{tabular}{|c|c|c|c|c|c|c|c|c|c|}
\hline
\(t\) & \(0\) & \(3\) & \(6\) & \(9\) & \(12\) & \(15\) & \(18\) & \(21\) & \(24\) \\
\hline
\(y\) & \(12\) & \(15.1\) & \(12.1\) & \(9.1\) & \(11.9\) & \(14.9\) & \(11.9\) & \(8.9\) & \(12.1\) \\
\hline
\end{tabular}
\end{center}

经长期观察，函数 \( y = f(t) \) 的图象可以近似地看成函数 \( y = k + A \sin(\omega t + \varphi) \) 的图象. 下面的函数中，最能近似表示表中数据间对应关系的函数是（\(\quad\)）

\choices
    {\( y = 12 + 3\sin \frac{\pi}{6} t \)，\(t \in [0, 24]\)}
    {\(y = 12 + 3\sin \left(\frac{\pi}{6} t + \pi\right)\)，\(t \in [0, 24]\)}
    {\(y = 12 + 3\sin \frac{\pi}{12} t\)，\(t \in [0, 24]\)}
    {\(y = 12 + 3\sin \left(\frac{\pi}{12} t + \frac{\pi}{2}\right)\)，\(t \in [0, 24]\)}
\end{problem}

\begin{answer}
A
\end{answer}

\begin{solution}
表中数据的平均水平约为 \(12\)，振幅约为 \(3\)，且在 \(t=0\) 时函数值约为中线，随后在 \(t=3\) 时达到最大值、在 \(t=9\) 时达到最小值，周期约为 \(12\). 选项 A 的周期为 \(\dfrac{2\pi}{\pi/6}=12\)，并且其在 \(t=0\)，\(3\)，\(6\)，\(9\)，\(12\) 时的函数值依次为 \(12\)，\(15\)，\(12\)，\(9\)，\(12\)，与表中数据最吻合. 选项 B 在 \(t=3\) 时取最小值，选项 C 的周期为 \(24\)，选项 D 在 \(t=0\) 时取值 \(15\)，均不符合，故选 A.
\end{solution}

\section{填空题}
\begin{problem}
设随机变量 \(\xi\) 的概率分布为 \(P(\xi = k) = \frac{a}{5^k}\)，\(a\) 为常数，\(k = 1\)，\(2\)，\(\dots\)，则 \(a = \)\fillinblank{}.
\end{problem}

\begin{answer}
\(4\)
\end{answer}

\begin{solution}
由概率总和为 \(1\)，得 \(\displaystyle\sum_{k=1}^{\infty}\frac{a}{5^k}=a\cdot\frac{1/5}{1-1/5}=\frac{a}{4}=1\)，所以 \(a=4\).
\end{solution}

\begin{problem}
将标号为 \(1\)，\(2\)，\(\cdots\)，\(10\) 的 \(10\) 个球放入标号为 \(1\)，\(2\)，\(\cdots\)，\(10\) 的 \(10\) 个盒子内，每个盒内放一个球，则恰好有 \(3\) 个球的标号与其所在盒子的标号不一致的放入方法共有\fillinblank{} 种. （以数字作答）
\end{problem}

\begin{answer}
\(240\)
\end{answer}

\begin{solution}
先选出标号不一致的 \(3\) 个盒子，共有 \(\mathrm C_{10}^{3}\) 种选法. 选定后，这 \(3\) 个球不能各自放回原盒，只能在这三个盒子中作无固定点排列；三个元素的错排有 \(2\) 种. 因此放入方法数为 \(\mathrm C_{10}^{3}\cdot2=120\cdot2=240\).
\end{solution}

\begin{problem}
设 \(A\)，\(B\) 为两个集合. 下列四个命题：

\begin{circlelist}
\item \(A \not\subseteq B \Leftrightarrow\) 对任意 \(x \in A\)，有 \(x \notin B\)；
\item \(A \not\subseteq B \Leftrightarrow A \cap B = \varnothing\)；
\item \(A \not\subseteq B \Leftrightarrow A \not\supseteq B\)；
\item \(A \not\subseteq B \Leftrightarrow\) 存在 \(x \in A\)，使得 \(x \notin B\).
\end{circlelist}

其中真命题的序号是\fillinblank{}. （把符合要求的命题序号都填上）
\end{problem}

\begin{answer}
\(4\)
\end{answer}

\begin{solution}
命题 1 把 \(A\not\subseteq B\) 错写成了对所有 \(x\in A\) 都有 \(x\notin B\)，而否定子集关系只需存在一个这样的元素，所以命题 1 为假. 命题 2 不成立，例如取 \(A=\varnothing\)，\(B=\varnothing\)，此时 \(A\cap B=\varnothing\)，但 \(A\subseteq B\). 命题 3 也不成立，例如取 \(A=\{1\}\)，\(B=\{1,2\}\)，则 \(A\subseteq B\)，但 \(A\not\supseteq B\). 命题 4 正是子集定义的否定形式：\(A\not\subseteq B\) 当且仅当存在 \(x\in A\) 使 \(x\notin B\). 故应填 \(4\).
\end{solution}

\begin{problem}
某日中午 \(12\) 时整，甲船自 \(A\) 处以 \(16\mathrm{~km/h}\) 的速度向正东行驶，乙船自 \(A\) 的正北 \(18\mathrm{~km}\) 处以 \(24\mathrm{~km/h}\) 的速度向正南行驶，则当日 \(12\) 时 \(30\) 分时两船之间距离对时间的变化率是\fillinblank{}\(\mathrm{km/h}\).
\end{problem}

\begin{answer}
\(-1.6\)
\end{answer}

\begin{solution}
设 \(t\) 为从中午 \(12\) 时起经过的小时数. 取 \(A\) 为原点，正东为 \(x\) 轴，则甲船坐标为 \((16t,0)\)，乙船坐标为 \((0,18-24t)\). 两船距离为 \(s(t)=\sqrt{(16t)^2+(18-24t)^2}\)，故
\(\displaystyle s'(t)=\frac{16^2t-24(18-24t)}{s(t)}\). 当 \(t=\dfrac12\) 时，\(s=\sqrt{8^2+6^2}=10\)，从而 \(\displaystyle s'\left(\frac12\right)=\frac{128-144}{10}=-1.6\mathrm{~km/h}\).
\end{solution}

\section{解答题}
\begin{problem}
已知 \(6\sin^2\alpha + \sin \alpha \cos \alpha - 2\cos^2\alpha = 0\)，\(\alpha \in \left[\frac{\pi}{2}, \pi\right]\)，求 \(\sin \left(2\alpha + \frac{\pi}{3}\right)\) 的值.
\end{problem}

\begin{solution}
由 \(6\sin^2\alpha+\sin\alpha\cos\alpha-2\cos^2\alpha=(3\sin\alpha+2\cos\alpha)(2\sin\alpha-\cos\alpha)\)，得 \(3\sin\alpha+2\cos\alpha=0\) 或 \(2\sin\alpha-\cos\alpha=0\). 在 \(\alpha\in[\frac\pi2,\pi]\) 时，\(\sin\alpha\geqslant0\)，\(\cos\alpha\leqslant0\)，第二个等式只能导致 \(\sin\alpha=\cos\alpha=0\)，不可能成立. 因此 \(3\sin\alpha+2\cos\alpha=0\). 又因该区间内实际只能有 \(\sin\alpha>0\)，故 \(\cos\alpha=-\dfrac32\sin\alpha\)，结合 \(\sin^2\alpha+\cos^2\alpha=1\) 得 \(\sin\alpha=\dfrac{2}{\sqrt{13}}\)，\(\cos\alpha=-\dfrac{3}{\sqrt{13}}\). 于是 \(\sin2\alpha=2\sin\alpha\cos\alpha=-\dfrac{12}{13}\)，\(\cos2\alpha=\cos^2\alpha-\sin^2\alpha=\dfrac{5}{13}\)，从而
\(\displaystyle\sin\left(2\alpha+\frac\pi3\right)=\sin2\alpha\cos\frac\pi3+\cos2\alpha\sin\frac\pi3=\frac{5\sqrt3-12}{26}\).
\end{solution}

\begin{problem}
如图，在棱长为 \(1\) 的正方体 \(ABCD - A_{1}B_{1}C_{1}D_{1}\) 中，点 \(E\) 是棱 \(BC\) 的中点，点 \(F\) 是棱 \(CD\) 上的动点.

\begin{enumerate}
\item 试确定点 \(F\) 的位置，使得 \(D_{1}E \perp\) 平面 \(AB_{1}F\)；
\item 当 \(D_{1}E \perp\) 平面 \(AB_{1}F\) 时，求二面角 \(C_{1} - EF - A\) 的大小. （结果用反三角函数值表示）
\end{enumerate}

\bitmapfigure[max width=0.44\linewidth]{img/2004/hubei_science/q18_fig1.png}
\end{problem}

\begin{solution}
\begin{enumerate}
\item 取 \(A\) 为原点，令 \(AB\)，\(AD\)，\(AA_1\) 分别为三个坐标轴的正方向. 设 \(DF=t\)，其中 \(0\leqslant t\leqslant1\)，则 \(E=(1,\frac12,0)\)，\(F=(t,1,0)\)，\(D_1=(0,1,1)\)，且可取 \(\overrightarrow{D_1E}=(1,-\frac12,-1)\)，\(\overrightarrow{AB_1}=(1,0,1)\)，\(\overrightarrow{AF}=(t,1,0)\). 有 \(\overrightarrow{D_1E}\cdot\overrightarrow{AB_1}=0\)，所以 \(D_1E\perp AB_1\)；又 \(\overrightarrow{D_1E}\cdot\overrightarrow{AF}=t-\frac12\). 因此 \(D_1E\perp AF\) 当且仅当 \(t=\frac12\)，从而 \(D_1E\perp\) 平面 \(AB_1F\) 当且仅当 \(F\) 是棱 \(CD\) 的中点.
\item 此时 \(E=(1,\frac12,0)\)，\(F=(\frac12,1,0)\)，所以 \(EF\parallel BD\). 设 \(H=AC\cap EF\)，则 \(H=(\frac34,\frac34,0)\)，且 \(AC\perp EF\). 因为 \(C_1H\) 在底面内的射影为 \(CH\)，也有 \(C_1H\perp EF\)，故二面角 \(C_1-EF-A\) 的平面角为 \(\angle AHC_1\). 在直角三角形 \(C_1CH\) 中，\(C_1C=1\)，\(CH=\frac{\sqrt2}{4}\)，所以 \(\tan\angle C_1HC=\dfrac{C_1C}{CH}=2\sqrt2\). 射线 \(HA\) 与 \(HC\) 反向，故 \(\angle AHC_1=\pi-\arctan(2\sqrt2)\).
\end{enumerate}
\end{solution}

\begin{problem}
如图，在 \(\mathrm{Rt}\triangle ABC\) 中，已知 \(BC = a\)，若长为 \(2a\) 的线段 \(PQ\) 以点 \(A\) 为中点，问 \(\overrightarrow{PQ}\) 与 \(\overrightarrow{BC}\) 的夹角 \(\theta\) 取何值时 \(\overrightarrow{BP} \cdot \overrightarrow{CQ}\) 的值最大？并求出这个最大值.

\bitmapfigure[float=inline,max width=0.30\linewidth]{img/2004/hubei_science/q19_fig1.png}
\FloatBarrier
\end{problem}

\begin{solution}
设 \(\bs{e}=\dfrac{1}{2a}\overrightarrow{PQ}\)，则 \(|\bs{e}|=1\). 由于 \(A\) 是 \(PQ\) 的中点，有 \(\overrightarrow{AP}=-a\bs{e}\)，\(\overrightarrow{AQ}=a\bs{e}\). 又 \(\triangle ABC\) 在 \(A\) 处为直角，所以 \(\overrightarrow{AB}\cdot\overrightarrow{AC}=0\). 因此
\(\displaystyle\overrightarrow{BP}\cdot\overrightarrow{CQ}=(-\overrightarrow{AB}-a\bs{e})\cdot(-\overrightarrow{AC}+a\bs{e})=-a^2+a\bs{e}\cdot(\overrightarrow{AC}-\overrightarrow{AB})=-a^2+a\bs{e}\cdot\overrightarrow{BC}\).
由于 \(|\overrightarrow{BC}|=a\)，且 \(\bs{e}\) 与 \(\overrightarrow{BC}\) 的夹角为 \(\theta\)，故 \(\overrightarrow{BP}\cdot\overrightarrow{CQ}=a^2(\cos\theta-1)\leqslant0\). 等号当且仅当 \(\cos\theta=1\)，即 \(\theta=0\)，所以最大值为 \(0\).
\end{solution}

\begin{problem}
直线 \(l: y = kx + 1\) 与双曲线 \(C: 2x^{2} - y^{2} = 1\) 的右支交于不同的两点 \(A\)，\(B\).

\begin{enumerate}
\item 求实数 \(k\) 的取值范围；
\item 是否存在实数 \(k\)，使得以线段 \(AB\) 为直径的圆经过双曲线 \(C\) 的右焦点 \(F\)？若存在，求出 \(k\) 的值. 若不存在，说明理由.
\end{enumerate}
\end{problem}

\begin{solution}
\begin{enumerate}
\item 联立直线与双曲线，得 \((k^2-2)x^2+2kx+2=0\). 要有两个不同交点，需判别式 \(\Delta=4(4-k^2)>0\)，且 \(k^2\ne2\). 又因为两交点在右支，两个根 \(x_1\)，\(x_2\) 均为正，所以 \(x_1x_2=\dfrac{2}{k^2-2}>0\)，从而 \(k^2>2\)；并且 \(x_1+x_2=\dfrac{2k}{2-k^2}>0\)，从而 \(k<0\). 综合得 \(-2<k<-\sqrt2\).
\item 双曲线的右焦点为 \(F=(c,0)\)，其中 \(c=\sqrt{\frac12+1}=\frac{\sqrt6}{2}\). 设 \(A=(x_1,kx_1+1)\)，\(B=(x_2,kx_2+1)\). 圆以 \(AB\) 为直径且经过 \(F\)，等价于 \(\overrightarrow{FA}\cdot\overrightarrow{FB}=0\)，即
\(\displaystyle(1+k^2)x_1x_2+(k-c)(x_1+x_2)+c^2+1=0\).
由韦达定理 \(x_1+x_2=\dfrac{2k}{2-k^2}\)，\(x_1x_2=\dfrac{2}{k^2-2}\)，代入并利用 \(c^2=\frac32\)，得 \(5k^2+4ck-6=0\)，即 \(5k^2+2\sqrt6 k-6=0\). 解得 \(k=-\dfrac{6+\sqrt6}{5}\) 或 \(k=\dfrac{6-\sqrt6}{5}\). 后者不属于第一问的范围，故存在且 \(k=-\dfrac{6+\sqrt6}{5}\).
\end{enumerate}
\end{solution}

\begin{problem}
某突发事件，在不采取任何预防措施的情况下发生的概率为 \(0.3\)；一旦发生，将造成 \(400\text{ 万元}\) 的损失. 现有甲，乙两种相互独立的预防措施可供采用. 单独采用甲，乙预防措施所需的费用分别为 \(45\text{ 万元}\) 和 \(30\text{ 万元}\). 采用相应预防措施后此突发事件不发生的概率分别是 \(0.9\) 和 \(0.85\). 若预防方案允许甲，乙两种预防措施单独采用，联合采用或不采用，请确定预防方案使总费用最少. 注：总费用 \(=\) 采取预防措施的费用 \(+\) 发生突发事件损失的期望值.
\end{problem}

\begin{solution}
不采取措施时，突发事件发生的概率为 \(0.3\)，总费用为 \(400\times0.3=120\text{ 万元}\). 单独采用甲时，事件发生概率为 \(1-0.9=0.1\)，损失期望值为 \(400\times0.1=40\text{ 万元}\)，总费用为 \(45+400\times0.1=85\text{ 万元}\). 单独采用乙时，事件发生概率为 \(1-0.85=0.15\)，损失期望值为 \(400\times0.15=60\text{ 万元}\)，总费用为 \(30+400\times0.15=90\text{ 万元}\). 联合采用时，由于两措施相互独立，事件发生概率为 \((1-0.9)(1-0.85)=0.015\)，总费用为 \(45+30+400\times0.015=81\text{ 万元}\). 比较四种方案，联合采用甲、乙两种预防措施时总费用最少，最少总费用为 \(81\text{ 万元}\).
\end{solution}

\begin{problem}
已知 \(a > 0\)，数列 \(\{a_n\}\) 满足 \(a_1 = a\)，\(a_{n+1} = a + \frac{1}{a_n}\)，\(n = 1\)，\(2\)，\(\ldots\).

\begin{enumerate}
\item 已知数列 \(\{a_n\}\) 极限存在且大于零，求 \(A = \displaystyle\lim_{n \to \infty} a_n\)（将 \(A\) 用 \(a\) 表示）；
\item 设 \(b_n = a_n - A\)，\(n = 1\)，\(2\)，\(\cdots\)，证明： \(b_{n+1} = -\frac{b_n}{A(b_n + A)}\).
\item 若 \(|b_n| \leqslant \frac{1}{2^n}\) 对 \(n = 1\)，\(2\)，\(\cdots\) 都成立，求 \(a\) 的取值范围.
\end{enumerate}
\end{problem}

\begin{solution}
\begin{enumerate}
\item 因为极限存在且 \(A>0\)，对递推式取极限得 \(A=a+\dfrac1A\)，即 \(A^2-aA-1=0\). 所以 \(A=\dfrac{a\pm\sqrt{a^2+4}}2\)，其中负号所得根小于零，故 \(\displaystyle A=\frac{a+\sqrt{a^2+4}}2\).
\item 由 \(A=a+\dfrac1A\) 及 \(a_n=b_n+A\)，有 \(\displaystyle b_{n+1}=a+\frac1{b_n+A}-A=-\frac1A+\frac1{b_n+A}=-\frac{b_n}{A(b_n+A)}\).
\item 由 \(b_1=a-A=-\dfrac1A\)，条件在 \(n=1\) 时给出 \(\dfrac1A\leqslant\dfrac12\)，即 \(A\geqslant2\). 又由 \(a=A-\dfrac1A\)，得 \(a\geqslant2-\dfrac12=\dfrac32\). 反之若 \(a\geqslant\dfrac32\)，则对应的正根 \(A\geqslant2\). 若 \(|b_n|\leqslant2^{-n}\)，则 \(b_n+A=a_n\geqslant A-|b_n|\geqslant\dfrac32\)，从而
\(\displaystyle |b_{n+1}|=\frac{|b_n|}{A|b_n+A|}\leqslant\frac{|b_n|}{3}\leqslant\frac{1}{3\cdot2^n}<\frac{1}{2^{n+1}}\). 由归纳法，所有 \(n\) 均满足所给不等式. 因此 \(a\) 的取值范围为 \(\left[\dfrac32,+\infty\right)\).
\end{enumerate}
\end{solution}
